The validated findings establish that robust improvements in the Prabhas\=a pipeline derive from two sources: architecture via RoPE, and pretraining objective design. RoPE applies complex-number rotations in embedding space to provide compositional structure across scales, consistent with widespread adoption in contemporary language model families such as Llama and Qwen. The pretraining objective exhibits marked scale dependence: at 10M parameters, pure MLM achieves 63.58 (plus or minus 1.73) versus hybrid MLM+CLM at 64.09 (plus or minus 0.26), with overlapping 95 percent confidence intervals indicating no significant difference. At 100M parameters, pure MLM yields 73.06 compared to 67.57 for hybrid MLM, a gain of 5.49 percentage points. This scale interaction reflects a fundamental trade-off whereby CLM's additional gradient signal may stabilize learning on small models where MLM signal is sparse, but becomes a bottleneck as scale increases. The Strict-Small submission (prabhasa-b\_ss, BLiMP 64.09) retains hybrid MLM+CLM, while the Strict track model (prabhasa-b\_s, BLiMP 73.06) adopts pure MLM to exploit larger scale. These scores position Prabhas\=a competitively relative to official baselines: 73.06 versus 74.53 on Strict and 64.09 versus 65.08 on Strict-Small, with differences within inter-seed variance.

In contrast, the k\=araka mechanisms tested do not carry measurable causal weight on standard linguistic diagnostics. Role-stratified masking (Arm K) yields negligible difference from uniform random masking (Arm C) at matched budget: on the targeted agreement and argument-structure paradigms, 82.03 versus 81.93 (95 percent CI minus 0.99 to plus 1.20), and on full BLiMP, 71.77 versus 70.49. Root-cause analysis revealed that earlier apparent gains were confounded by mask-rate variation and frequency weighting, not role distribution itself. A supervised auxiliary objective for k\=araka role prediction shows a similar null pattern: the five-seed mean difference is 0.76 percentage points (95 percent CI minus 0.74 to plus 2.27, not statistically significant), attenuating monotonically from 2.65 points at seed one to minus 0.34 at seed five. A specificity control using shuffled-role labels demonstrates that the auxiliary effect is not attributable to generic regularization: the shuffled variant yields 0.58 benefit versus 0.76 for genuine roles. These null results may reflect scope limitations of the 10M-100M parameter regime, single-language training, and evaluation restricted to BLiMP agreement tasks rather than compositional generalization benchmarks such as COGS or SCAN, which directly test argument-structure learning (see Appendices \ref{app:training} and \ref{app:results}).

The k\=araka masking and auxiliary objective findings do not falsify the linguistic hypothesis that argument structure matters for pretraining; rather, they narrow scope by demonstrating that the specific mechanisms tested do not yield measurable advantage at this scale and evaluation surface. Larger models, longer training runs, multilingual settings, or evaluation on compositional generalization tasks may exhibit different behavior. A single-seed run at 100M provides no basis for ruling out such effects at 350M scale or above. The pre-registered, budget-matched comparison provides a valid and conservative closure under stated conditions.

Prabhasa demonstrates a methodological contribution to rigorous pretraining research. Pre-registration of three hypotheses with effect-size thresholds enforces discipline and prevents outcome-driven narrative construction. Matched-budget comparisons for the k\=araka mechanism studies eliminate confounds that could inflate efficacy estimates. Multi-seed validation in F1 (three seeds at 10M) and F3 (five seeds at 10M) exposes inter-seed variance that single-seed publications often obscure. Reporting null results with full transparency, including seed-level values, confidence intervals, and specificity controls, contributes to research culture in which null findings are validated outcomes rather than failures. In the small-model regime where computational resources are constrained, clarity about what does and does not work provides value for downstream researchers.

Several limitations circumscribe the scope. The 10M to 100M parameter range is orders of magnitude smaller than contemporary pretraining; whether null findings on k\=araka mechanisms generalize to larger models remains unknown. Evaluation on BLiMP agreement and argument-structure diagnostics may not be optimal for structure-aware pretraining; compositional generalization benchmarks could reveal benefits that BLiMP does not capture. The k\=araka role assignments derive from BPE-token-level heuristics rather than full syntactic parsing, introducing noise; high-accuracy parses would clarify whether nulls reflect role-assignment error or genuine effects. The Strict-track submission model (100M, single seed) is underpowered for confident claims; a three-to-five seed replication would provide tighter bounds.

Future work should pursue four directions. First, systematic characterization of the objective-scale interaction: at what model size does pure MLM surpass hybrid approaches, and does the crossover point vary across BLiMP, SCAN, and COGS? Second, targeted evaluation of k\=araka mechanisms on compositional generalization benchmarks (COGS, SCAN, CFQ) that directly test argument-structure learning. Third, multilingual extension using Sanskrit as a bridge language, testing whether structure-aware mechanisms show stronger signatures in morphologically rich source languages. Fourth, evaluation of synergy between the best H1 arm (pure-MLM at 100M) and planned post-training enhancements (the Navya-Ny\=aya reasoning scaffold) on inference-quality metrics. The validated findings on architecture and objective choice, combined with transparent null results on k\=araka mechanisms and grounded in pre-registration and matched-budget discipline, establish a foundation for future work across larger scales and alternative evaluation surfaces.
